\clearpage
\VerdierPhysicalPageAnchor{0145}%
\begingroup
\small
\setlength{\parskip}{0pt}%
\setlength{\emergencystretch}{1.5em}%
\setlength{\abovedisplayskip}{3pt plus 1pt minus 1pt}%
\setlength{\belowdisplayskip}{3pt plus 1pt minus 1pt}%

\VerdierSemanticAnchor{ch2-2-3-2-free-object-continuation}
\noindent 以 $C$ 与 $C'$ 为元素的宇宙后表述为：若函子
$Z\longmapsto\operatorname{Hom}_C(Z,Y)$ 与
$Z\longmapsto\operatorname{Hom}_{C'}(F(Z),X)$ 同构，则 $Y$ 在 $X$ 上
右自由。

\VerdierSemanticAnchor{ch2-2-3-2-right-left-terminology}
若存在 $C'$ 的对象 $X$ 使 $Y$ 在 $X$ 上右自由，则称 $C$ 的对象 $Y$
\emph{$F$-右自由}。若存在 $C$ 的对象 $Y$ 在 $X$ 上右自由，则称
$C'$ 的对象 $X$ \emph{$F$-右可自由化}。通过取反范畴，以同样方式定义
\emph{左}自由对象与左可自由化对象。

\VerdierSemanticAnchor{ch2-2-3-3}
\medskip\begingroup\itshape
\noindent\textbf{命题 2.3.3.}\quad
设 $\mathcal D$ 为三角范畴，$\mathcal B$ 为满三角子范畴，
$Q\colon\mathcal D\longrightarrow\mathcal D/\mathcal B$ 为商函子
(2.2.10)，$S_{\mathcal D}(\mathcal B)$ 为 $\mathcal D$ 中与
$\mathcal B$ 关联的乘法系 (2.1.8)。

\VerdierSemanticAnchor{ch2-2-3-3-a}
\par\noindent a) 对 $\mathcal D$ 的每个对象 $Y$，下列条件等价：
\VerdierSemanticAnchor{ch2-2-3-3-a-i}
\par\noindent\hspace{1.2em}i) 对象 $Y$ 是 $Q$-右自由的。
\VerdierSemanticAnchor{ch2-2-3-3-a-ii}
\par\noindent\hspace{1.2em}ii) 对每个态射
$s\in S_{\mathcal D}(\mathcal B)$，$s\colon S\longrightarrow T$，映射
\[
  \operatorname{Hom}_{\mathcal D}(s,Y)\colon
  \operatorname{Hom}_{\mathcal D}(T,Y)\longrightarrow
  \operatorname{Hom}_{\mathcal D}(S,Y)
\]
是双射。
\VerdierSemanticAnchor{ch2-2-3-3-a-iii}
\par\noindent\hspace{1.2em}iii) 对 $\mathcal B$ 的每个对象 $B$，
$\operatorname{Hom}_{\mathcal D}(B,Y)=0$。
\VerdierSemanticAnchor{ch2-2-3-3-a-iv}
\par\noindent\hspace{1.2em}iv) 每个态射 $Y\xrightarrow{s}Z$，其中
$s\in S_{\mathcal D}(\mathcal B)$，都有收缩。
\VerdierSemanticAnchor{ch2-2-3-3-a-v}
\par\noindent\hspace{1.2em}v) 对 $\mathcal D$ 的每个对象 $X$，映射
\[
  \operatorname{Hom}_{\mathcal D}(X,Y)\longrightarrow
  \operatorname{Hom}_{\mathcal D/\mathcal B}(Q(X),Q(Y))
\]
是双射。

\VerdierSemanticAnchor{ch2-2-3-3-b}
\par\noindent b) 设 $L_R(Q)$ 为由 $Q$-右自由对象定义的
$\mathcal D$ 的满子范畴。范畴 $L_R(Q)$ 是三角、严格满且饱和的。
\VerdierSemanticAnchor{ch2-2-3-3-c}
\par\noindent c) 商函子在 $L_R(Q)$ 上的限制
\[
  Q|_{L_R(Q)}\colon L_R(Q)\longrightarrow\mathcal D/\mathcal B
  ,
\]
充分忠实，并且在对象上单射。
\VerdierSemanticAnchor{ch2-2-3-3-d}
\par\noindent d) 设
$\operatorname{Lib}_R(Q)\subset\mathcal D/\mathcal B$ 为由
$\mathcal D/\mathcal B$ 中的 $Q$-右可自由化对象定义的满子范畴。
$\operatorname{Lib}_R(Q)$ 是严格满三角子范畴。若可数积或可数直和在
$\mathcal D$ 中可表示，则 $\operatorname{Lib}_R(Q)$ 是饱和的。
\VerdierSemanticAnchor{ch2-2-3-3-e}
\par\noindent e) 函子 $Q$ 把 $L_R(Q)$ 映入 $\operatorname{Lib}_R(Q)$。
$\operatorname{Lib}_R(Q)$ 是 $L_R(Q)$ 的像的严格满闭包。自由箭头都是同构。
\endgroup
\endgroup
